\documentclass[10pt]{article}
\usepackage[utf8]{inputenc}
\usepackage[preprint]{tmlr}
\usepackage{hyperref}
\usepackage{url}
\usepackage{amsmath,amssymb}
\usepackage{longtable,booktabs,array,calc}
\providecommand{\tightlist}{\setlength{\itemsep}{0pt}\setlength{\parskip}{0pt}}
\providecommand{\pandocbounded}[1]{#1}
\DeclareUnicodeCharacter{00A7}{\S{}}
\DeclareUnicodeCharacter{00BD}{\ensuremath{\tfrac{1}{2}}}
\DeclareUnicodeCharacter{2153}{\ensuremath{\tfrac{1}{3}}}
\DeclareUnicodeCharacter{00D7}{\ensuremath{\times}}
\DeclareUnicodeCharacter{03A3}{\ensuremath{\Sigma}}
\DeclareUnicodeCharacter{03B2}{\ensuremath{\beta}}
\DeclareUnicodeCharacter{03B4}{\ensuremath{\delta}}
\DeclareUnicodeCharacter{03B5}{\ensuremath{\varepsilon}}
\DeclareUnicodeCharacter{03B7}{\ensuremath{\eta}}
\DeclareUnicodeCharacter{03B9}{\ensuremath{\iota}}
\DeclareUnicodeCharacter{2013}{--}
\DeclareUnicodeCharacter{2014}{---}
\DeclareUnicodeCharacter{2016}{\ensuremath{\|}}
\DeclareUnicodeCharacter{2026}{\ldots{}}
\DeclareUnicodeCharacter{2212}{\ensuremath{-}}
\DeclareUnicodeCharacter{2265}{\ensuremath{\geq}}
\DeclareUnicodeCharacter{00B9}{\textsuperscript{1}}
\DeclareUnicodeCharacter{00B3}{\textsuperscript{3}}
\DeclareUnicodeCharacter{2074}{\textsuperscript{4}}
\DeclareUnicodeCharacter{2075}{\textsuperscript{5}}
\DeclareUnicodeCharacter{207B}{\textsuperscript{-}}
\DeclareUnicodeCharacter{2082}{\textsubscript{2}}
\DeclareUnicodeCharacter{2084}{\textsubscript{4}}
\DeclareUnicodeCharacter{2032}{\ensuremath{'}}
\DeclareUnicodeCharacter{2154}{\ensuremath{\tfrac{2}{3}}}
\DeclareUnicodeCharacter{2081}{\textsubscript{1}}
\DeclareUnicodeCharacter{2208}{\ensuremath{\in}}
\DeclareUnicodeCharacter{03BB}{\ensuremath{\lambda}}
\DeclareUnicodeCharacter{03BC}{\ensuremath{\mu}}
\DeclareUnicodeCharacter{03B3}{\ensuremath{\gamma}}
\DeclareUnicodeCharacter{03B1}{\ensuremath{\alpha}}
\DeclareUnicodeCharacter{2264}{\ensuremath{\leq}}
\title{An exact world audits its own experiments: what a small stochastic kernel can and cannot settle about learners}
\author{\name W. Anthony Mason \email fsgeek@cs.ubc.ca \\ \addr University of British Columbia}
\def\month{MM}
\def\year{YYYY}
\def\openreview{\url{https://openreview.net/forum?id=XXXX}}
\begin{document}
\maketitle
\begin{abstract}
Exact synthetic worlds let us compare what a learned model acquires with what the data could teach it. We build one
in which the observation channel is the controlled variable: a small stochastic operating-system kernel, observed
through five nested interfaces, with exact posteriors and predictive laws for every window an observer can see. The
uncertainty an interface leaves about the next event splits exactly into randomness that no observation could remove
and hidden state the interface withholds. The mean randomness is the same on every interface by construction; the
withheld part falls from about half of the uncertainty on the sparsest interface to about a third on the richest. We
then use the same exact quantities to audit which learner experiments the world can support. The natural contrast
between pure randomness and withheld state, at matched uncertainty, has about 1 \% of the evaluated population on its
withheld side, and 86--95 \% of its random side is a single mechanism. Whether one prediction shows where its
uncertainty comes from depends on the question and the precision. Among exact laws, telling pure randomness from
substantial withheld state becomes ambiguous between separations of 0.01 and 0.02 in total variation, while the
amount of randomness stays readable to 0.02 and becomes ambiguous by 0.05. For twelve trained transformers, measured by
the prior mass of the classes whose exact law is at least as close to the model's output as the truth's, the kind of
uncertainty meets a pre-declared readability rule in every cell and the amount of randomness in none. On the next
event's kind, a comparison that separates an input rewrite from its training target suggests that the target
contributes to the rewrite's gain; whether the input adds anything is unresolved at three seeds. An exact oracle makes
such limitations measurable; it does not guarantee that the desired contrast exists.

\end{abstract}
\hypertarget{introduction}{%
\section{Introduction}\label{introduction}}

\textbf{Exact worlds are popular for a reason.} A transformer trained on Othello moves can be probed against the true board
\citep{li2022emergent}, and transformers trained on sequences from hidden Markov models
\citep{shai2024belief, agarwal2025bayesian} or on automata with known entropies \citep{someya2025locality} can be compared with
exact Bayes quantities, because the information in the training data is known. A shortfall from the exact ceiling is
then a property of the learner, not of unknown data.

\textbf{Exactness removes one excuse, not all of them.} \citet{vafa2024world} show that in exact finite-automaton worlds a standard
diagnostic can still overstate how much of the world a model has recovered. We report a complementary set of
limitations from one exact world, each measured with the world's own exact quantities: whether a planned contrast has
support in the world, at what precision a question about a model's outputs becomes answerable, and whether an
information-preserving intervention changes only what it was meant to change.

\textbf{The world.} Yupana is a small stochastic kernel (four threads, two CPUs, two locks, an asynchronous device) whose
hidden state and transition law we own. The hidden system is held fixed while the \emph{observation interface} varies: five
nested interfaces each show one more field of every event record. For every window an observer can see, we compute
the exact posterior over hidden state and the exact law of the next event.

\textbf{Contributions.}

\begin{enumerate}
\def\labelenumi{\arabic{enumi}.}
\tightlist
\item
  \textbf{Exact measurements of what each interface makes knowable} (§3). Next-event uncertainty splits into randomness and
  withheld state; by construction, interfaces can change only the second on average. On the content interfaces, state uncertainty
  overstates the predictively relevant part by a factor of 3--5 while ranking fields identically. Where in the window
  the knowable information sits differs by interface, and the profiles cross.
\item
  \textbf{A support audit} (§4). For a planned contrast: the eligible mass, the matched mass, and the mechanism
  composition of the pure-randomness side, all computable from the exact tables before any training (we ran them retrospectively); and a
  precision sweep for an output-level identifiability question, placed beside the measured precision of trained
  models.
\item
  \textbf{A worked learner case} (§5), including a comparison that separates an intervention's input from its training
  target.
\end{enumerate}

\textbf{Motivation.} The instrument was motivated by an impossibility result about verifying a model's epistemic state from
its text alone \citep{epistemic2026observability}: in a world where the truth is known, one can measure how much the observation
channel withholds and how much a model's outputs reveal. §4.3 is the nearest this paper comes to that question, and it
answers it only for one-step predictions in this world. It is an analogy to that result, not a test of it.

\hypertarget{the-instrument}{%
\section{The instrument}\label{the-instrument}}

\hypertarget{the-world}{%
\subsection{The world}\label{the-world}}

Yupana is a small stochastic operating-system kernel: four threads, two CPUs, two locks and one asynchronous device,
with a request pool of four and a device queue of depth two. Each thread runs a fixed looping program, giving four
contention roles. The kernel makes two kinds of random choice. The scheduler picks the next thread round-robin with
probability 1 − ε and uniformly among runnable threads with probability ε; we report ε = 1 and ε = ½. The device
completes a pending request on a step with probability 1/3, the request chosen uniformly among those pending. Every
choice is labelled in the next record by its kind and actor, so the future stays random however much is observed.

An episode runs 40 records from reset. The observer sees a \emph{window}: the last 8 records ending at an endpoint drawn
uniformly from \{2, 4, \ldots, 40\}, with a RESET token when the window reaches back to the reset. This window law was fixed
before any learner was trained.

\hypertarget{interfaces}{%
\subsection{Interfaces}\label{interfaces}}

Every record has five fields: KIND, ACTOR, OBJECT (which lock or device), RELATED (the thread woken, or the owner) and
LINEAGE (the request identifier). Interface r0 shows only KIND, and each of r1--r4 shows one more field. Hidden fields
are rendered as a placeholder, so every record is five tokens on every interface. Before any learner existed, each added
field was shown to change some posterior or some finite-horizon predictive distribution.

Entity names carry no identity across episodes: they are drawn per episode from pools much larger than the entity
counts (64 thread names for four threads, for example).

\hypertarget{exact-quantities-and-their-validation}{%
\subsection{Exact quantities and their validation}\label{exact-quantities-and-their-validation}}

Two observers appear. The \emph{structural observer} sees structural identities; we use it only in §3.2. The \emph{pattern
observer} sees each entity renamed by order of first appearance, so that only the pattern of who recurs remains.
Because names are redrawn every episode, it is the learner's comparison point, and every other quantity in this paper
is for it, at position 7 (after seven records, predicting the eighth). The \emph{ceiling} is the law-mass-weighted entropy
of the exact predictive law given the pattern class: a Bayes predictive loss.

\begin{longtable}[]{@{}lll@{}}
\toprule
\begin{minipage}[b]{0.30\columnwidth}\raggedright
quantity\strut
\end{minipage} & \begin{minipage}[b]{0.30\columnwidth}\raggedright
validation\strut
\end{minipage} & \begin{minipage}[b]{0.30\columnwidth}\raggedright
coverage\strut
\end{minipage}\tabularnewline
\midrule
\endhead
\begin{minipage}[t]{0.30\columnwidth}\raggedright
structural window posteriors\strut
\end{minipage} & \begin{minipage}[t]{0.30\columnwidth}\raggedright
the filtering recursion agrees with an independent path summation\strut
\end{minipage} & \begin{minipage}[t]{0.30\columnwidth}\raggedright
exhaustive: 8,542,940 windows, all ten (interface, ε) cells\strut
\end{minipage}\tabularnewline
\begin{minipage}[t]{0.30\columnwidth}\raggedright
pattern-class predictive laws\strut
\end{minipage} & \begin{minipage}[t]{0.30\columnwidth}\raggedright
agree exactly with an independent pattern filter\strut
\end{minipage} & \begin{minipage}[t]{0.30\columnwidth}\raggedright
exhaustive at a small window law (episodes of 8 records, windows of 4, endpoint stride 2); at the frozen law, a random sample of 82--91 test windows per cell, all exactly equal\strut
\end{minipage}\tabularnewline
\begin{minipage}[t]{0.30\columnwidth}\raggedright
lag profile ι\_m = H\_m − H*\strut
\end{minipage} & \begin{minipage}[t]{0.30\columnwidth}\raggedright
the KL route and the entropy route agree\strut
\end{minipage} & \begin{minipage}[t]{0.30\columnwidth}\raggedright
every (interface, ε, m), to 3 × 10⁻¹⁴\strut
\end{minipage}\tabularnewline
\begin{minipage}[t]{0.30\columnwidth}\raggedright
ceiling-holding rewrites\strut
\end{minipage} & \begin{minipage}[t]{0.30\columnwidth}\raggedright
equal KIND marginal per class; equal multiset of record probabilities per class; an independent construction checked by brute force\strut
\end{minipage} & \begin{minipage}[t]{0.30\columnwidth}\raggedright
exhaustive over classes\strut
\end{minipage}\tabularnewline
\begin{minipage}[t]{0.30\columnwidth}\raggedright
the R/G split (§3.1)\strut
\end{minipage} & \begin{minipage}[t]{0.30\columnwidth}\raggedright
each interface's mean H* reproduces the ceiling; G ≥ −10⁻¹⁵ per class (a mutual information, up to rounding)\strut
\end{minipage} & \begin{minipage}[t]{0.30\columnwidth}\raggedright
every class, both ε\strut
\end{minipage}\tabularnewline
\bottomrule
\end{longtable}

\textbf{What none of this validates is the kernel.} Every route calls the same kernel functions, so the gates check the
computation \emph{given} the kernel. The kernel rests on its tests, its reviews, and one property tested directly: from a
reset, the r1 projection determines the structural state for every transition path up to eight records, at both ε.
Appendix B specifies the kernel and the exact computation.

\hypertarget{learners}{%
\subsection{Learners}\label{learners}}

A decoder-only transformer (4 layers, width 128, 838,940 parameters) is trained from scratch on 2 M windows per cell
with AdamW and a cosine schedule, for 5 or 15 epochs. At 15 epochs validation loss is still falling slowly, so we write
``at 15 epochs'' rather than ``at convergence''. The measurement is the KIND \emph{shortfall} at position 7: the realised
log-ratio between the exact law and the model's prediction, whose class mean estimates KL(P* ‖ q). On the content
interfaces it runs at about 1--4 \% of the ceiling. Here q is the model's distribution renormalised only to remove the
end-of-window token, so probability placed on non-KIND tokens is penalised; §4.3 instead renormalises q over the KINDs,
because it compares distributions over KIND. The trained model is the unit of inference: three seeds per cell, and
inference per model or per matched seed pair. Appendix A lists which experiment was pre-registered, on which models and
at what budget.

\hypertarget{what-each-interface-lets-an-observer-know-exactly}{%
\section{What each interface lets an observer know, exactly}\label{what-each-interface-lets-an-observer-know-exactly}}

Unless stated otherwise, the observer is the pattern observer, the position is after seven records, and the target is
the next record's KIND. Position-7 quantities cover the windows that run at least seven records: 0.85 of window law
mass.

\hypertarget{where-the-uncertainty-comes-from}{%
\subsection{Where the uncertainty comes from}\label{where-the-uncertainty-comes-from}}

The next-KIND entropy H*(c) of a class c splits exactly into two parts:

\begin{itemize}
\tightlist
\item
  \textbf{R(c) = Σ\_s P(s \textbar{} c) H(KIND \textbar{} s)}, the randomness an observer who knew the hidden state would still face;
\item
  \textbf{G(c) = H*(c) − R(c) = I(KIND; S \textbar{} c)}, the part due to state the interface withholds.
\end{itemize}

\begin{longtable}[]{@{}llllll@{}}
\toprule
ε = 1 & r0 & r1 & r2 & r3 & r4\tabularnewline
\midrule
\endhead
H*, the ceiling & 1.176 & 0.999 & 0.955 & 0.911 & 0.911\tabularnewline
R, transition randomness & 0.609 & 0.609 & 0.609 & 0.609 & 0.609\tabularnewline
G, observation gap & 0.567 & 0.391 & 0.346 & 0.302 & 0.302\tabularnewline
\bottomrule
\end{longtable}

\emph{Bits, law-mass means. At ε = ½: R = 0.544 on every interface; G = 0.560, 0.329, 0.279, 0.240, 0.240.}

R's mean is the same on every interface \textbf{by construction}: it is E{[}H(KIND \textbar{} S){]}, which no observation changes. Per
class, R does vary. The finding is in G: about half of the sparsest interface's mean next-KIND uncertainty is withheld
state, and about a third of r3's (r4 agrees with r3 to within 10⁻⁵ bits on this target). On average, added fields act
only on the withheld part; the mean randomness is beyond any interface's reach.

\hypertarget{knowing-the-state-is-not-predicting-the-next-event}{%
\subsection{Knowing the state is not predicting the next event}\label{knowing-the-state-is-not-predicting-the-next-event}}

The classical figure of merit for telemetry is identifiability: how much uncertainty about the hidden state the record
leaves. Log placement work minimises a quantity of this kind \citep{zhao2017log20}. For the structural observer at the
window's last record, on the whole-record target:

\begin{longtable}[]{@{}llllll@{}}
\toprule
ε = 1 & r0 & r1 & r2 & r3 & r4\tabularnewline
\midrule
\endhead
state entropy H(S \textbar{} window) & 3.947 & 0.601 & 0.528 & 0.405 & 0.390\tabularnewline
next-record observation gap & 0.442 & 0.150 & 0.135 & 0.106 & 0.103\tabularnewline
\bottomrule
\end{longtable}

Most state uncertainty is irrelevant to the next record: on r1, 0.60 bits of state uncertainty correspond to 0.15 bits
of reducible next-record uncertainty. (The gap is a mutual information with the state, so it cannot exceed the state
entropy; the size of the difference is the finding.) On the content interfaces, each further field (OBJECT, RELATED,
LINEAGE) removes 3--5 times more of the first than of the second (4.7, 4.3, 4.4 at ε = 1; 4.3, 3.1, 4.0 at ε = ½), and
both measures rank those fields the same way. The first field, ACTOR, is not comparable, because it also changes the
next-record target (r0's record is KIND alone); its ratio is 11.5 at ε = 1. Along this nested ladder, an
identifiability objective would choose the same fields as a predictor's ceiling while overstating the relevant
uncertainty by a fairly stable factor; that need not hold for unrestricted field selection. This table uses a different
observer, position and target from §3.1, so the two should not be compared number by number.

\hypertarget{where-in-the-window-the-knowable-information-sits}{%
\subsection{Where in the window the knowable information sits}\label{where-in-the-window-the-knowable-information-sits}}

For a class c, let ι\_m(c) be the information about the next KIND that lies more than m records back: the KL divergence
from the full-window law to the law given only the last m records. Its law-mass mean is the interface's \emph{lag profile}.
It equals the order-m entropy minus the ceiling: a per-class, exact relative of the m-local entropy of
\citet{someya2025locality}, with two differences of convention (our order m counts records and conditions on m of them, where
theirs conditions on m − 1 symbols), computed for each interface of one fixed system.

The profiles cross. At lags 1--3, r3 and r4 (effectively tied; r4 is larger by under 10⁻⁵ bits) leave the most
information beyond m records; from m = 4, r2 does (ι₄ = 0.282 against r3's 0.259, ε = 1). A richer field both adds
information and moves information that was already available closer to the point of use, and the ceiling registers
only the first. The crossing is about 0.02 bits in one world; we report it as an illustration of why a profile is
needed, not as a general effect.

\hypertarget{rewrites-that-hold-the-ceiling}{%
\subsection{Rewrites that hold the ceiling}\label{rewrites-that-hold-the-ceiling}}

A \emph{window-level} rewrite may fill a masked slot with any function of the window's earlier records. It adds nothing an
observer of the window lacked, so the ceiling is unchanged class by class, which three gates verify exactly (equal KIND
marginals per class; equal multisets of record probabilities per class; an independent construction checked by brute
force on small worlds). We use one such rewrite in §5: the \emph{relay}, which writes into each record the KIND of the same
actor's previous record in the window.

\hypertarget{which-contrasts-can-this-world-carry-and-at-what-precision}{%
\section{Which contrasts can this world carry, and at what precision?}\label{which-contrasts-can-this-world-carry-and-at-what-precision}}

Exact ceilings tell a learner study what is knowable. They do not tell it whether the comparison it wants exists in the
world: a contrast needs classes on both sides that match on everything but the quantity of interest. Nor do they tell
it whether a question about a model's outputs can be answered at the precision those outputs have. The first check,
and the exact half of the second, need only the exact class tables, so they could run before any model is trained;
measuring a model's own precision needs the model. We ran them afterwards, on the learner studies of §5, and report what
they would have shown.

\hypertarget{a-motivating-error}{%
\subsection{A motivating error}\label{a-motivating-error}}

Our first reading of the §5.1 results took a trained model's shortfall to be a \emph{step}: roughly constant once a class has
any far information (ι₄ \textgreater{} 0), rather than rising with it. That reading came from fits that controlled the class's
next-KIND entropy, and it does not survive. Classes with ι₄ = 0 take only three next-KIND entropy values, so the entropy
control was extrapolating far outside the support of one side of the contrast. In raw per-model bins, shortfall rises
with ι₄ above ι₄ = 0.001 in every model. Counting the entropy values on each side of ι₄ = 0 would have shown the problem
before the fit.

\hypertarget{randomness-against-withheld-state}{%
\subsection{Randomness against withheld state}\label{randomness-against-withheld-state}}

§3.1 splits each class's next-KIND uncertainty exactly into randomness R and withheld state G. The natural learner
question is whether a model treats ``the world is random here'' differently from ``the telemetry hides the state here''.
Its cleanest form compares classes of matched H* whose uncertainty is pure randomness (G = 0, H* \textgreater{} 0) with classes
where much of it is withheld (G/H* ≥ 0.5), matched in 0.05-bit H* bins.

\begin{longtable}[]{@{}lllllll@{}}
\toprule
\begin{minipage}[b]{0.12\columnwidth}\raggedright
ε\strut
\end{minipage} & \begin{minipage}[b]{0.12\columnwidth}\raggedright
interface\strut
\end{minipage} & \begin{minipage}[b]{0.12\columnwidth}\raggedright
pure: eligible and matched (\%)\strut
\end{minipage} & \begin{minipage}[b]{0.12\columnwidth}\raggedright
of which the device coin (\%)\strut
\end{minipage} & \begin{minipage}[b]{0.12\columnwidth}\raggedright
high-gap: eligible (\%)\strut
\end{minipage} & \begin{minipage}[b]{0.12\columnwidth}\raggedright
high-gap: matched (\%)\strut
\end{minipage} & \begin{minipage}[b]{0.12\columnwidth}\raggedright
graded variation (\%)\strut
\end{minipage}\tabularnewline
\midrule
\endhead
\begin{minipage}[t]{0.12\columnwidth}\raggedright
1\strut
\end{minipage} & \begin{minipage}[t]{0.12\columnwidth}\raggedright
r1 / r2 / r3\strut
\end{minipage} & \begin{minipage}[t]{0.12\columnwidth}\raggedright
14.6 / 15.6 / 16.7\strut
\end{minipage} & \begin{minipage}[t]{0.12\columnwidth}\raggedright
92 / 91 / 86\strut
\end{minipage} & \begin{minipage}[t]{0.12\columnwidth}\raggedright
39 / 35 / 28\strut
\end{minipage} & \begin{minipage}[t]{0.12\columnwidth}\raggedright
1.4 / 1.2 / 0.9\strut
\end{minipage} & \begin{minipage}[t]{0.12\columnwidth}\raggedright
21 / 24 / 19\strut
\end{minipage}\tabularnewline
\begin{minipage}[t]{0.12\columnwidth}\raggedright
½\strut
\end{minipage} & \begin{minipage}[t]{0.12\columnwidth}\raggedright
r1 / r2 / r3\strut
\end{minipage} & \begin{minipage}[t]{0.12\columnwidth}\raggedright
16.1 / 17.2 / 18.3\strut
\end{minipage} & \begin{minipage}[t]{0.12\columnwidth}\raggedright
95 / 93 / 89\strut
\end{minipage} & \begin{minipage}[t]{0.12\columnwidth}\raggedright
36 / 32 / 26\strut
\end{minipage} & \begin{minipage}[t]{0.12\columnwidth}\raggedright
1.4 / 1.3 / 0.8\strut
\end{minipage} & \begin{minipage}[t]{0.12\columnwidth}\raggedright
39 / 44 / 38\strut
\end{minipage}\tabularnewline
\bottomrule
\end{longtable}

\emph{Shares of the position-7 population (law mass 0.85; r4 agrees with r3 at this precision). Pure: G = 0, H* \textgreater{} 0;
high-gap: G/H* ≥ 0.5; matched: sharing a bin floor(H*/0.05) with the other side. Every pure class is matched, and the
matched mass lies in three shared bins. Device coin: a pure class whose H* is the entropy of a 1/3 coin (0.918 bits). Every such class has next-KIND law
⅓ IO\_COMPLETE and ⅔ on a single other event: the kernel gives probability 1/3 to a completion whenever a request is
pending and the remaining 2/3 to the scheduling step, which in these classes has one outcome (Appendix B).
Graded variation: mass in bins where G/H* spans 0.30 or more between its mass-weighted 10th and 90th percentiles;
restricting to bins with at least 0.5 \% of mass and five classes gives 17--22 \% (ε = 1) and 36--43 \% (ε = ½).}

Two things follow. The clean contrast is narrow: its high-gap side is about 1 \% of the population, and 86--95 \% of its
pure-randomness side is one mechanism (a pending I/O completes with probability 1/3). A graded contrast has far more
mass, but we have not audited how many mechanisms it spans and do not claim it is balanced. Small mass does not by itself
make a contrast unusable; it says how much of the evaluation the contrast can draw on, and from what. The narrowness is
a consequence of design: the kernel places randomness only in the scheduler and the device, to keep exact filtering
inside its support budget.

\hypertarget{can-one-prediction-say-where-the-uncertainty-comes-from}{%
\subsection{Can one prediction say where the uncertainty comes from?}\label{can-one-prediction-say-where-the-uncertainty-comes-from}}

In general, the split between randomness and withheld state is not identifiable from a predictive distribution, even an
exact one \citep{bergna2026decoupled}; that result concerns unknown generative models. With the kernel known, structure can
restore identification, so we ask the question inside the world for a stated family: the kernel, the interface and ε are
known, and the observer is the pattern observer at position 7. Two senses of ``where it comes from'' matter, and they
behave differently:

\begin{itemize}
\tightlist
\item
  the \emph{kind}: pure randomness (G = 0 with H* \textgreater{} 0) against substantial withheld state (G ≥ 0.1 bits); classes in
  between have no kind. This differs from §4.2's high-gap side (G/H* ≥ 0.5), which asks for classes where withheld
  state dominates, for a matched learner contrast; here the question is only whether any substantial withheld state
  is present;
\item
  the \emph{amount of randomness} R.
\end{itemize}

\textbf{Exact laws.} Classes with identical exact next-KIND laws almost always share their split. The exceptions are below
0.05 \% of position-7 mass on every interface, but real: on r3 at ε = 1, 24 reachable classes share the law ⅓ ACQUIRE,
⅓ IO\_COMPLETE, ⅓ RELEASE, with splits including (R, G) = (0.918, 0.667) and (1.585, 0).

\textbf{Laws a small distance apart.} Let δ be a total-variation separation between two classes' exact laws. Two laws can both
lie within η of one prediction only if they are within 2η of each other, so δ corresponds to a prediction error of about
δ/2. For each class we ask whether some class within δ differs from it. These are existential statistics: one rare
neighbour suffices.

\begin{longtable}[]{@{}lll@{}}
\toprule
\begin{minipage}[b]{0.30\columnwidth}\raggedright
δ\strut
\end{minipage} & \begin{minipage}[b]{0.30\columnwidth}\raggedright
kind differs (ε = 1 / ½, r1--r4)\strut
\end{minipage} & \begin{minipage}[b]{0.30\columnwidth}\raggedright
R differs by ≥ 0.1 bits (ε = 1, r1 and r3)\strut
\end{minipage}\tabularnewline
\midrule
\endhead
\begin{minipage}[t]{0.30\columnwidth}\raggedright
0.01\strut
\end{minipage} & \begin{minipage}[t]{0.30\columnwidth}\raggedright
0.1--0.2 \% / 0.6--1.7 \%\strut
\end{minipage} & \begin{minipage}[t]{0.30\columnwidth}\raggedright
---\strut
\end{minipage}\tabularnewline
\begin{minipage}[t]{0.30\columnwidth}\raggedright
0.02\strut
\end{minipage} & \begin{minipage}[t]{0.30\columnwidth}\raggedright
14--17 \% / 17--19 \%\strut
\end{minipage} & \begin{minipage}[t]{0.30\columnwidth}\raggedright
1.8--2.4 \%\strut
\end{minipage}\tabularnewline
\begin{minipage}[t]{0.30\columnwidth}\raggedright
0.05\strut
\end{minipage} & \begin{minipage}[t]{0.30\columnwidth}\raggedright
16--18 \% / 20--23 \%\strut
\end{minipage} & \begin{minipage}[t]{0.30\columnwidth}\raggedright
43--44 \%\strut
\end{minipage}\tabularnewline
\bottomrule
\end{longtable}

\emph{Shares of all position-7 mass. Opposite-sign R/G swaps of at least 0.05 bits each occur on about 3 \% of mass at
δ = 0.02 and 18--21 \% at δ = 0.05 (ε = 1, r1 and r3).}

The two senses fail differently over the separations evaluated. The kind becomes ambiguous between δ = 0.01 and 0.02
and changes little from there to δ = 0.05, the largest separation we evaluated. R stays pinned up to δ = 0.02 and is
ambiguous on 43--44 \% of mass by 0.05. (G itself, the amount of withheld state, differs
on 49--51 \% of mass at δ = 0.02, mostly because a small perturbation moves H*.)

\textbf{Measured on trained models.} We ran twelve trained models (r1 and r3, both ε, three seeds each) on their test
windows. For each window, the \emph{look-alike set} is the set of classes whose exact law is at least as close to the model's
output as the true class's law is: \(A(w) = \{c' : \mathrm{TV}(P^*_{c'}, q_w) \le \mathrm{TV}(P^*_{c}, q_w)\}\), where \(c\) is the window's class and \(q_w\) the model's output. The
statistic is a(w) = Σ m(c′) over the c′ in A(w) that differ from c, divided by Σ m(c′) over all of A(w), with m a
class's law mass. The set is calibrated by an oracle that knows the true class, so these are prior-mass shares of
look-alikes, not error rates of any monitor. Before each run we declared a result \emph{readable} if the mean of a(w) is
below 0.05 and fewer than 10 \% of windows have a(w) ≥ 0.10, \emph{not readable} if the mean is at least 0.20 or at least
40 \% of windows have a(w) ≥ 0.10, and otherwise neither; all three seeds of a cell must agree.

\begin{itemize}
\tightlist
\item
  The models' TV distance to the exact law: median 0.015--0.027, mean 0.027--0.041; 19--31 \% of windows exceed 0.05.
\item
  Look-alike mass of the other \emph{kind}: 0.5--1.4 \% on average; 1.3--3.9 \% of windows have 10 \% or more.
\item
  Look-alike mass differing in R by 0.1 bits or more: 5--11 \% on average; 10--23 \% of windows have 10 \% or more, higher at
  ε = ½ and on r1 than on r3.
\end{itemize}

Three details of the kind measurement dilute its shares. Every class with G = 0, including deterministic ones, counted
as pure; windows whose class has no kind (16--18 \% of mass) score zero; and a window's own R and G were taken as the mean
over classes sharing its exact law, which differs from the class's own value only on collision groups (below 0.05 \% of
mass). Deterministic classes and classes with no kind together carry 32--39 \% of the mass, so at most 39 \% of windows score
zero by construction, and shares conditioned on the rest are at most 1/0.61
times those above: under 2.3 \% on average and 6.4 \% of windows.

We declared rules for calling each result readable or not before each run. The R measurement came first and fell
between the rules in every cell. The kind measurement was declared after it, at a reviewer's prompting, and before it
ran; it meets the ``readable'' rule in every cell and seed, with or without the correction above. For R we report the
numbers only. Set against the exact sweep, the models' median error (0.015--0.027) allows separations of up to roughly 0.03--0.05,
where the kind is already ambiguous on about a fifth of the mass existentially; that is only a compatibility bound.
The look-alike measurement asks the more specific question at the models' actual outputs, and there the look-alike mass
of the other kind stays small: ambiguous neighbours exist, but they are rare.

\hypertarget{the-audit-as-a-procedure}{%
\subsection{The audit as a procedure}\label{the-audit-as-a-procedure}}

The checks use only the exact class tables, which a study computing exact ceilings already has:

\begin{enumerate}
\def\labelenumi{\arabic{enumi}.}
\tightlist
\item
  \textbf{State the estimand and the matching variables} (e.g.~the shortfall difference between G = 0 and G/H* ≥ 0.5,
  matched on H* to 0.05 bits).
\item
  \textbf{Compute the eligible and matched mass on each side}, per interface and ε, as shares of the population the learner
  will be evaluated on.
\item
  \textbf{Report the mechanism composition} of each side, for example by grouping classes by their exact next-event law or by
  the kernel rule that produces it, with the dominant share.
\item
  \textbf{State the admissible family} for any identifiability claim (what a monitor is assumed to know), and test it inside
  the world, distinguishing the senses of the question.
\item
  \textbf{Sweep precision}: repeat any identifiability check over a range of separations; once models exist, place their
  measured output precision on the same axis.
\end{enumerate}

The exact parts run on a workstation CPU from the existing tables. The model measurement needs one forward pass per
model over its test windows; each twelve-model run took under a quarter of an hour of wall time as parallel CPU
processes.

\hypertarget{the-learner-case}{%
\section{The learner case}\label{the-learner-case}}

The learners and the shortfall measure are those of §2.4. This is a measured case, not a general account of learners.

\hypertarget{shortfall-is-higher-where-far-information-is}{%
\subsection{Shortfall is higher where far information is}\label{shortfall-is-higher-where-far-information-is}}

Within each content interface, a pre-registered regression of per-window KIND shortfall d(w) on the class's far
information, d(w) = α + β ι₄(c) + γ (ι₁(c) − ι₄(c)) + λ H*(c) + μ log₂ m(c) + e(w), with near information ι₁ − ι₄,
ceiling entropy H* and log law mass m, gives β \textgreater{} 0 with t ≥ 3 in every run:
21 of 21 held-out runs at 5 epochs and 24 of 24 at 15 epochs. It survives errors clustered by class, and fixed effects
for the last record's kind crossed with the number of distinct actors (lowest t 3.44), in all 45 of those evaluations
(runs, not independent seeds). An exploratory check at about four times the parameters, in one cell, keeps the sign. The
same analysis for the actor field passes its pre-registered test but is not robust to the fixed effects (9 of 24 runs
keep t ≥ 3 at 15 epochs). The association is a correlation within cells: classes with more far information may also be
rarer or harder in ways the controls miss. A pre-registered prediction that the coefficient would shrink no faster than
overall shortfall with training failed in 8 of 8 cells. The size of the association at 15 epochs:

\begin{longtable}[]{@{}lllll@{}}
\toprule
\begin{minipage}[b]{0.17\columnwidth}\raggedright
ε\strut
\end{minipage} & \begin{minipage}[b]{0.17\columnwidth}\raggedright
interface\strut
\end{minipage} & \begin{minipage}[b]{0.17\columnwidth}\raggedright
mean shortfall, ι₄ = 0\strut
\end{minipage} & \begin{minipage}[b]{0.17\columnwidth}\raggedright
mean shortfall, ι₄ \textgreater{} 0.4\strut
\end{minipage} & \begin{minipage}[b]{0.17\columnwidth}\raggedright
β with fixed effects (t)\strut
\end{minipage}\tabularnewline
\midrule
\endhead
\begin{minipage}[t]{0.17\columnwidth}\raggedright
1\strut
\end{minipage} & \begin{minipage}[t]{0.17\columnwidth}\raggedright
r1 / r2 / r3 / r4\strut
\end{minipage} & \begin{minipage}[t]{0.17\columnwidth}\raggedright
0.001--0.004\strut
\end{minipage} & \begin{minipage}[t]{0.17\columnwidth}\raggedright
0.023--0.050\strut
\end{minipage} & \begin{minipage}[t]{0.17\columnwidth}\raggedright
0.009--0.023 (3.6--8.1)\strut
\end{minipage}\tabularnewline
\begin{minipage}[t]{0.17\columnwidth}\raggedright
½\strut
\end{minipage} & \begin{minipage}[t]{0.17\columnwidth}\raggedright
r1 / r2 / r3 / r4\strut
\end{minipage} & \begin{minipage}[t]{0.17\columnwidth}\raggedright
0.001--0.004\strut
\end{minipage} & \begin{minipage}[t]{0.17\columnwidth}\raggedright
0.035--0.061\strut
\end{minipage} & \begin{minipage}[t]{0.17\columnwidth}\raggedright
0.013--0.024 (3.6--9.6)\strut
\end{minipage}\tabularnewline
\bottomrule
\end{longtable}

\emph{Bits; ranges over the three seeds of each interface. β is bits of shortfall per bit of far information.}

\textbf{Baselines.} A Witten--Bell count table over the ceiling's own classes, trained on the same windows, has 1.33--5.22 times
the KIND shortfall of every transformer on every content interface. A parameter-matched GRU on the transformer's untuned recipe
was pre-registered with a floor fixed before the models ran: each had to beat the count table before its association
could be read. All eight failed it (KIND 0.045--0.076 bits against the table's 0.039--0.064), so their large, positive
association statistics are uninformative under the declared rule.

\hypertarget{a-copy-placed-at-the-point-of-use-separating-input-from-supervision}{%
\subsection{A copy placed at the point of use: separating input from supervision}\label{a-copy-placed-at-the-point-of-use-separating-input-from-supervision}}

The relay rewrite (§3.4) holds the ceiling identical class by class while putting the same actor's previous KIND in each
record. All three relay models had lower KIND shortfall than all three controls. A pre-registered prediction that the gain
would scale with how much far information the relay moves closer failed. In the original training recipe the copy is
both an input and a training target. Two pre-registered follow-ups separated the two by varying only the OBJECT slot.
The five arms below are the cells that exist; no arm shows the copy while training on a placeholder, so the design is a
set of controlled contrasts, not a full factorial (KIND and ACTOR shortfall at 15 epochs, seeds 1--3):

\begin{longtable}[]{@{}lllll@{}}
\toprule
arm & OBJECT input & OBJECT target & KIND & ACTOR\tabularnewline
\midrule
\endhead
r1 & placeholder & placeholder & 0.0157, 0.0146, 0.0158 & 0.0312, 0.0299, 0.0314\tabularnewline
r1m & placeholder & none & 0.0176, 0.0209, 0.0170 & 0.0346, 0.0438, 0.0409\tabularnewline
r1Rs & placeholder & relay value & 0.0125, 0.0141, 0.0142 & 0.0296, 0.0400, 0.0342\tabularnewline
r1Rm & relay copy & none & 0.0281, 0.0411, 0.0307 & 0.0328, 0.0363, 0.0345\tabularnewline
r1R & relay copy & relay value & 0.0130, 0.0095, 0.0125 & 0.0278, 0.0262, 0.0308\tabularnewline
\bottomrule
\end{longtable}

A contrast is \emph{resolved} when every model of one arm lies beyond every model of the other in the pre-registered
direction. With three seeds a side, exact paired sign-flip p-values cannot fall below 1/8, and none is conventionally
significant. On KIND:

\begin{itemize}
\tightlist
\item
  \textbf{An unscored copy hurt} (r1Rm against r1m): showing the copy without training on it cost the learner, relative to a
  constant placeholder. This cannot distinguish an informative unscored copy from any varying field.
\item
  \textbf{Changing the slot's target moved the shortfall} (r1Rs against r1): training the slot to produce the relay value,
  without ever showing it, lowered KIND shortfall in every seed, with a margin of 0.0005 bits between the closest models
  (unrounded values). Removing the slot from the loss altogether (r1m against r1) raised KIND shortfall by a comparable
  amount. KIND shortfall thus moves with any change to the OBJECT slot's loss, so r1Rs's gain cannot be credited to the
  content of the lookup rather than to having a nontrivial target in that slot. (Leaving OBJECT out of the loss also shrinks the loss denominator by about 24 \%.)
\item
  \textbf{Whether the copy adds to supervision is unresolved} (r1R against r1Rs).
\end{itemize}

On ACTOR, none of these contrasts resolves. The simple placement hypothesis, ``a near copy helps a learned reader'', finds
no support here: the original gain does not establish an input-placement benefit. The training target contributes, and
the input's increment remains open at three seeds.

For the audit of §4, this is a second kind of limitation an exact world leaves in place. Holding the ceiling fixed
controls what is \emph{knowable}. It does not control what the learner is \emph{trained to produce}, and an information-preserving
rewrite of the input can change the training targets with it.

\hypertarget{pitfalls-we-met-that-are-not-specific-to-exact-worlds}{%
\section{Pitfalls we met that are not specific to exact worlds}\label{pitfalls-we-met-that-are-not-specific-to-exact-worlds}}

Several of our errors are standard, and we only list them. Between-seed variance dominated paired comparisons, so the
trained model is the unit \citep{bouthillier2021variance, henderson2017matters}. Far information measured on the original
history and on the relay-rewritten history (§3.4) correlate at 0.79 over the paired test windows; in an exploratory
follow-up with a declared reading, a comparison of which one tracks learner shortfall separated them, and in the one cell
then checked the separation vanished once a reviewer added fixed effects for the last record's kind. Competing exact
explanations should be compared under simple observable controls. Some output quantities are fixed by the training
objective under stated conditions: a model whose log loss is stationary with respect to a global temperature, on the
distribution and positions where it is evaluated, has zero mean one-step entropy-calibration gap
\citep{braverman2019calibration}, and KL-regularised post-training has closed-form optima \citep{rafailov2023dpo}. Such quantities
are reference predictions, not findings. And baselines need a quality floor fixed before they run (§5.1).

\hypertarget{related-work}{%
\section{Related work}\label{related-work}}

\textbf{Exact testbeds.} Othello-GPT \citep{li2022emergent} and its successors; belief-state geometry in transformers trained on
HMM data \citep{shai2024belief} and in pretrained LLMs given HMM data in context \citep{balcells2026belief}; Bayesian ``wind
tunnels'' \citep{agarwal2025bayesian} and an entropy-aligned axis in production LLMs \citep{agarwal2025scaling}. These comparisons
vary the process (sampled transition and emission matrices) or the model; ours holds one process fixed and varies its
observation interface. \citet{vafa2024world} evaluate world models in exact automaton worlds and show that common diagnostics
mislead, and \citet{vafa2025found} probe foundation models for world models through their inductive bias, in settings
including orbital mechanics. Our audit is a different check: support and precision, not recovery.

\textbf{Locality and entropy convergence.} \citet{someya2025locality} relate learning to m-local entropy on natural and synthetic
languages, measure learner excess loss against exact PFSA entropies, and use entropy-preserving transformations as
interventions. Our lag profile is a per-class relative of their quantity (counted in records, conditioning on m of
them), and our ceiling-preserving rewrites are close to their interventions; what differs is the fixed hidden system,
class-by-class preservation, and the separation of input from supervision. Block-entropy convergence
\citep{crutchfield2003regularities} is the classical form.

\textbf{Separating kinds of uncertainty.} The split of predictive uncertainty into aleatoric and epistemic parts is not
identifiable from a predictive distribution in general \citep{bergna2026decoupled}, whose decomposition is of predictive
variance; ours is of entropy, conditioned on the hidden state. Joint or paired predictions can identify it
\citep{osband2021epistemic, johnson2024experts}. The closest precedent is \citet[Section 5.1]{ahdritz2024knowable}: in a synthetic task
whose questions carry an explicit bit marking each answer as fixed or resampled, one-step predictions are near 50/50 in
both cases, and the difference appears only in how the model updates on a provided answer. Our §4.3 asks the one-step
question in a world with no such cue, whose kernel is known, and finds that the answer depends on which sense of the
question is asked.

\textbf{Overlap.} The support audit is positivity checking \citep{petersen2012positivity, damour2017overlap} that can be applied
before a learner study, using exact class masses.

\textbf{Telemetry design.} Log placement work optimises the identifiability of execution from logs \citep{zhao2017log20}; §3.2
computes that kind of quantity beside the predictive one.

\hypertarget{limits}{%
\section{Limits}\label{limits}}

One world, with randomness placed by design in two mechanisms; one window law; one architecture (with an exploratory
check at about four times the parameters in one cell); three seeds per cell; KIND at position 7. The identifiability
results hold for the stated admissible family (known kernel, interface and ε, pattern observer, position 7). Learner
precision was measured on twelve models on two interfaces. The kernel itself is validated only by its tests, its
reviews and one directly tested property, and the pattern-class laws at the frozen law are validated on a sample of
82--91 test windows per cell rather than exhaustively. A reparameterised world, varying completion and dispatch probabilities with
the posterior support unchanged, would test whether the audit's findings belong to this world's mechanisms.

\section*{Contribution statement}
This research was conducted by successive instances of Claude (Anthropic) that owned the project, each continuing from
the last through a timestamped Git history and written handoffs. The founding proposal was drafted with a ChatGPT
(OpenAI) instance. Designs, pre-registrations and results were reviewed adversarially by separate Claude and Codex
(OpenAI) instances; the repository records, for each, whether its review preceded the run or followed it. The reviews
contributed substance as well as corrections, including the thesis sentence that closes the abstract and the collision
example of §4.3. W. Anthony Mason originated the project, provided the hardware and the research framework (the
\emph{ayllu}, a collaboration organised around reciprocity between human and AI participants), took part as a collaborator,
and posed questions that redirected the work; he is the author of record on behalf of the collaboration. Summary
results and review records are in the repository; per-class and per-window data are regenerable from it and pinned by
hash.

\bibliography{paper1m}
\bibliographystyle{tmlr}
\appendix
\hypertarget{which-models-carry-which-claim}{%
\section{Which models carry which claim}\label{which-models-carry-which-claim}}

\begin{longtable}[]{@{}lllll@{}}
\toprule
\begin{minipage}[b]{0.17\columnwidth}\raggedright
experiment\strut
\end{minipage} & \begin{minipage}[b]{0.17\columnwidth}\raggedright
cells\strut
\end{minipage} & \begin{minipage}[b]{0.17\columnwidth}\raggedright
seeds per cell\strut
\end{minipage} & \begin{minipage}[b]{0.17\columnwidth}\raggedright
budget\strut
\end{minipage} & \begin{minipage}[b]{0.17\columnwidth}\raggedright
analysis\strut
\end{minipage}\tabularnewline
\midrule
\endhead
\begin{minipage}[t]{0.17\columnwidth}\raggedright
far-information association, KIND (held-out)\strut
\end{minipage} & \begin{minipage}[t]{0.17\columnwidth}\raggedright
5 of the 8 r1--r4 cells\strut
\end{minipage} & \begin{minipage}[t]{0.17\columnwidth}\raggedright
3; 6 in two cells\strut
\end{minipage} & \begin{minipage}[t]{0.17\columnwidth}\raggedright
5 epochs\strut
\end{minipage} & \begin{minipage}[t]{0.17\columnwidth}\raggedright
pre-registered before evaluation (21 runs)\strut
\end{minipage}\tabularnewline
\begin{minipage}[t]{0.17\columnwidth}\raggedright
far-information association, ACTOR\strut
\end{minipage} & \begin{minipage}[t]{0.17\columnwidth}\raggedright
8 cells\strut
\end{minipage} & \begin{minipage}[t]{0.17\columnwidth}\raggedright
36 runs: 21 held-out, 15 existing\strut
\end{minipage} & \begin{minipage}[t]{0.17\columnwidth}\raggedright
5 epochs\strut
\end{minipage} & \begin{minipage}[t]{0.17\columnwidth}\raggedright
pre-registered before any actor quantity\strut
\end{minipage}\tabularnewline
\begin{minipage}[t]{0.17\columnwidth}\raggedright
15-epoch ladder\strut
\end{minipage} & \begin{minipage}[t]{0.17\columnwidth}\raggedright
8 r1--r4 cells, plus r0 at both ε\strut
\end{minipage} & \begin{minipage}[t]{0.17\columnwidth}\raggedright
3\strut
\end{minipage} & \begin{minipage}[t]{0.17\columnwidth}\raggedright
15 epochs\strut
\end{minipage} & \begin{minipage}[t]{0.17\columnwidth}\raggedright
pre-registered before evaluation; three r1 ε = 1 models seen earlier\strut
\end{minipage}\tabularnewline
\begin{minipage}[t]{0.17\columnwidth}\raggedright
relay\strut
\end{minipage} & \begin{minipage}[t]{0.17\columnwidth}\raggedright
r1 against r1R, ε = 1\strut
\end{minipage} & \begin{minipage}[t]{0.17\columnwidth}\raggedright
3 per arm\strut
\end{minipage} & \begin{minipage}[t]{0.17\columnwidth}\raggedright
5 and 15 epochs\strut
\end{minipage} & \begin{minipage}[t]{0.17\columnwidth}\raggedright
pre-registered prediction (failed)\strut
\end{minipage}\tabularnewline
\begin{minipage}[t]{0.17\columnwidth}\raggedright
loss-masked arm\strut
\end{minipage} & \begin{minipage}[t]{0.17\columnwidth}\raggedright
r1Rm\strut
\end{minipage} & \begin{minipage}[t]{0.17\columnwidth}\raggedright
3\strut
\end{minipage} & \begin{minipage}[t]{0.17\columnwidth}\raggedright
15 epochs\strut
\end{minipage} & \begin{minipage}[t]{0.17\columnwidth}\raggedright
pre-registered\strut
\end{minipage}\tabularnewline
\begin{minipage}[t]{0.17\columnwidth}\raggedright
relay follow-up contrasts\strut
\end{minipage} & \begin{minipage}[t]{0.17\columnwidth}\raggedright
r1m, r1Rs\strut
\end{minipage} & \begin{minipage}[t]{0.17\columnwidth}\raggedright
3 per arm\strut
\end{minipage} & \begin{minipage}[t]{0.17\columnwidth}\raggedright
15 epochs\strut
\end{minipage} & \begin{minipage}[t]{0.17\columnwidth}\raggedright
pre-registered contrasts\strut
\end{minipage}\tabularnewline
\begin{minipage}[t]{0.17\columnwidth}\raggedright
count baseline\strut
\end{minipage} & \begin{minipage}[t]{0.17\columnwidth}\raggedright
10 cells\strut
\end{minipage} & \begin{minipage}[t]{0.17\columnwidth}\raggedright
---\strut
\end{minipage} & \begin{minipage}[t]{0.17\columnwidth}\raggedright
---\strut
\end{minipage} & \begin{minipage}[t]{0.17\columnwidth}\raggedright
exploratory\strut
\end{minipage}\tabularnewline
\begin{minipage}[t]{0.17\columnwidth}\raggedright
GRU baseline\strut
\end{minipage} & \begin{minipage}[t]{0.17\columnwidth}\raggedright
8 r1--r4 cells\strut
\end{minipage} & \begin{minipage}[t]{0.17\columnwidth}\raggedright
1\strut
\end{minipage} & \begin{minipage}[t]{0.17\columnwidth}\raggedright
15 epochs\strut
\end{minipage} & \begin{minipage}[t]{0.17\columnwidth}\raggedright
pre-registered, with a quality floor\strut
\end{minipage}\tabularnewline
\begin{minipage}[t]{0.17\columnwidth}\raggedright
larger model\strut
\end{minipage} & \begin{minipage}[t]{0.17\columnwidth}\raggedright
r1 ε = 1\strut
\end{minipage} & \begin{minipage}[t]{0.17\columnwidth}\raggedright
3 (15 epochs), 1 (45 epochs)\strut
\end{minipage} & \begin{minipage}[t]{0.17\columnwidth}\raggedright
15, 45 epochs\strut
\end{minipage} & \begin{minipage}[t]{0.17\columnwidth}\raggedright
exploratory\strut
\end{minipage}\tabularnewline
\begin{minipage}[t]{0.17\columnwidth}\raggedright
learner precision (§4.3)\strut
\end{minipage} & \begin{minipage}[t]{0.17\columnwidth}\raggedright
r1, r3 at both ε\strut
\end{minipage} & \begin{minipage}[t]{0.17\columnwidth}\raggedright
3\strut
\end{minipage} & \begin{minipage}[t]{0.17\columnwidth}\raggedright
15 epochs\strut
\end{minipage} & \begin{minipage}[t]{0.17\columnwidth}\raggedright
two declared measurements; the second declared after the first was read\strut
\end{minipage}\tabularnewline
\bottomrule
\end{longtable}

\textbf{Training.} AdamW (β 0.9, 0.95; weight decay 0.1), a 500-step linear warm-up to a learning rate of 10⁻³, cosine decay
to 10⁻⁴, batch 256, gradient clipping at 1.0; the loss is the mean cross-entropy over scored target positions. At 15
epochs the last two epochs account for a median of 7.5 \% of the decline in validation loss from epoch 5 (at most 17 \%)
across the 30 ladder runs. Test windows come from episodes disjoint from training.

\textbf{Scoring.} For the realised next record with KIND x, the KIND shortfall is log₂ P*(x \textbar{} class) − log₂ \(\tilde{q}\)(x), where \(\tilde{q}\) is
the model's softmax at the KIND position renormalised to remove only the end-of-window token. Class means estimate
KL(P* ‖ \(\tilde{q}\)); single windows can be negative. Subsets are always chosen by functions of the prefix, never by another
realised field.

\hypertarget{the-kernel-and-the-exact-computation}{%
\section{The kernel and the exact computation}\label{the-kernel-and-the-exact-computation}}

\textbf{State.} For each of the four threads: a program counter and a status (RUNNABLE; RUNNING; blocked on lock l; blocked
on I/O request r; blocked on a full device queue; TERMINATED). Also: the set of threads on the two CPUs, each lock's
owner and FIFO wait queue, the device queue of (thread, request) pairs (capacity two), and a round-robin cursor (absent
from the dynamics at ε = 1). An episode starts with every thread RUNNABLE at program counter 0, no thread running, both
locks free and the device queue empty.

\textbf{Programs} (each loops forever): thread 0 acquires lock 0, releases it, issues I/O; thread 1 acquires lock 0 then
lock 1 and releases them in reverse order; thread 2 computes, then acquires and releases lock 0; thread 3 issues I/O,
then acquires and releases lock 1.

\textbf{One step emits one record.} If a request is pending, the device completes one of the pending requests, chosen
uniformly, with total probability 1/3 (record IO\_COMPLETE, with the issuer as actor and the request as lineage); the
issuer wakes. The remaining probability goes to the scheduling stage. If a CPU is free and a thread is runnable, a
DISPATCH places a thread on a CPU for one instruction (uniform over runnable threads with probability ε, round-robin
with probability 1 − ε). Otherwise one running thread, chosen by the same rule, executes its next instruction: STEP for computation; ACQUIRE, or
BLOCK with the owner as RELATED if the lock is held; RELEASE, with the lock handed directly to the head of its wait
queue, who appears as RELATED; IO\_ISSUE with a fresh request identifier as LINEAGE, or BLOCK on the device if its queue is full.
If nothing can happen, the step is IDLE. A thread leaves the CPU after each instruction. The record kinds are thus
DISPATCH, STEP, ACQUIRE, BLOCK, RELEASE, IO\_ISSUE, IO\_COMPLETE and IDLE, plus a RESET token at the start of a window that
reaches back to the reset.

\textbf{The exact computation.} One recursion per (window law, ε) runs at the richest interface; coarser interfaces are
projections. Its frontier maps (state, last-8-record window) to exact probability: at each tick every (state, window)
is pushed through the kernel's transition distribution, and at each endpoint the frontier's mass is added, divided by
the number of endpoints, to the joint law of (window, state). A \emph{pattern class} at position k is a window prefix of k
records after renaming each entity type by order of first appearance, together with whether the window starts at a
reset. Its law mass is the summed mass of its structural windows, and its next-record law is the posterior over
(state, binding of pattern names to structural entities), pushed through one kernel step and re-keyed through the
binding. Every probability is an exact rational; logarithms are taken in floating point at the end. Equality checks
between independent routes are exact equality of rationals.

\textbf{Size and cost} at the frozen law: reachable states peak at 1,953 (ε = 1) and 4,173 (ε = ½); the recursion's frontier
peaks at 1,964,919 (state, window) pairs; one pass takes about nine minutes and about 5 GB on one workstation CPU.
Pattern classes at position 7 number 5,540 (r0), 89,474 (r1), 116,624 (r2), 146,455 (r3) and 160,363 (r4).

\end{document}
